%%%%%%%%%%%%%%%%%%%%%%%%%%%%%%%%%%%%%%%%%%%%%%%
%
% Template per Elaborato di Laurea
% DISI - Dipartimento di Ingegneria e Scienza dell’Informazione
%
% update 2015-09-10
%
% Per la generazione corretta del 
% pdflatex nome_file.tex
% bibtex nome_file.aux
% pdflatex nome_file.tex
% pdflatex nome_file.tex
%
%%%%%%%%%%%%%%%%%%%%%%%%%%%%%%%%%%%%%%%%%%%%%%%

% formato FRONTE RETRO
\documentclass[a4paper]{book}
\usepackage{epsfig}
\usepackage{lipsum}
\usepackage{graphicx}
\usepackage{pgffor}
\usepackage{epigraph}
%\usepackage{plain}
%\usepackage{setspace}
% \usepackage[paperheight=29.7cm,paperwidth=21cm,outer=1.5cm,inner=2.5cm,top=2cm,bottom=2cm]{geometry} % per definizione layout
% \usepackage{titlesec} % per formato custom dei titoli dei capitoli
\usepackage[hidelinks]{hyperref}
\hypersetup{pdftex}

%%%%%%%%%%%%%%
% supporto lettere accentate
%
% unnecessary: using XeLaTeX
%
%\usepackage[latin1]{inputenc} % per Windows;
%\usepackage[utf8x]{inputenc} % per Linux (richiede il pacchetto unicode);
%\usepackage[applemac]{inputenc} % per Mac.

%\singlespacing

\usepackage[english]{babel}
\usepackage{csquotes}

\def\mtoday{\the\year-\twodigits\month-\twodigits\day}
\def\twodigits#1{\ifnum#1<10 0\fi\the#1}

\usepackage[backend=biber, dateabbrev=false, urldate=comp]{biblatex}
\addbibresource{biblio.bib}

\usepackage{xcolor}
\usepackage{listings}
\usepackage{amsmath}
\usepackage{amsthm}

\newtheorem{proposition}{Proposition}
\newtheorem{lemma}{Lemma}
\newtheorem{corollary}{Corollary}
\theoremstyle{definition}
\newtheorem{definition}{Definition}
\theoremstyle{remark}
\newtheorem{axiom}{Assumption}
\newtheorem{observation}{Observation}
\newtheorem{example}{Example}

\newcommand{\source}[1]{\ensuremath{\text{source}\left[{#1}\right]}}
\newcommand{\dest}[1]{\ensuremath{\text{dest}\left[{#1}\right]}}


\begin{document}

%    \pagestyle{plain}

    % nessuna numerazione
    \pagenumbering{gobble} 
    \thispagestyle{empty}

\begin{center}
  \begin{figure}[h!]
    \centering{\psfig{file=logo_unitn_black_center.eps}}
%    \centering{\includegraphics{newlogo_unitn_it, width=0.6\linewidth}}
  \end{figure}
  

  \vspace{1 cm} 

  \LARGE{Department of Information Engineering and Computer Science\\}

  \vspace{1 cm} 
  \Large{Bachelor's Degree in \\
    Computer Science
    %Ingegneria dell'Informazione e delle Comunicazioni
    %Ingegneria dell'Informazione e Organizzazione d'Impresa
    %Ingegneria Elettronica e delle Telecomunicazioni
  }

  \vspace{2 cm} 
  \Large\textsc{Final Dissertation\\} 
  \vspace{1 cm} 
  \Huge\textsc{Titolo\\}
  \Large{\it{Sottotitolo (alcune volte lungo - opzionale)}}


  \vspace{2 cm} 
  \begin{tabular*}{\textwidth}{ c @{\extracolsep{\fill}} c }
  \Large{Supervisor} & \Large{Student}\\
  \Large{Prof. Gabriel Mark Kuper}& \Large{Daniele Parmeggiani}\\
  \end{tabular*}

  \vspace{2 cm} 

  \Large{Academic year 2020/2021}
  
\end{center}


    \cleardoublepage
 
%%%%%%%%%%%%%%%%%%%%%%%%%%%%%%%%%%%%%%%%%%%%%%%%%%%%%%%%%%%%%%%%%%%%%%%%%%
%%%%%%%%%%%%%%%%%%%%%%%%%%%%%%%%%%%%%%%%%%%%%%%%%%%%%%%%%%%%%%%%%%%%%%%%%%
%% Nota
%%%%%%%%%%%%%%%%%%%%%%%%%%%%%%%%%%%%%%%%%%%%%%%%%%%%%%%%%%%%%%%%%%%%%%%%%%
%% Sezione Ringraziamenti opzionale
%%%%%%%%%%%%%%%%%%%%%%%%%%%%%%%%%%%%%%%%%%%%%%%%%%%%%%%%%%%%%%%%%%%%%%%%%%
%%%%%%%%%%%%%%%%%%%%%%%%%%%%%%%%%%%%%%%%%%%%%%%%%%%%%%%%%%%%%%%%%%%%%%%%%%
%  \emph{
%    To my father, and my mother who supported me throughout these years.
%  }
    \begin{center}
    \emph{Dedication}
    \end{center}
    \clearpage
%  \pagestyle{plain} % nessuna intestazione e pie pagina con numero al centro

  
    % inizio numerazione pagine in numeri arabi
    \mainmatter

%%%%%%%%%%%%%%%%%%%%%%%%%%%%%%%%%%%%%%%%%%%%%%%%%%%%%%%%%%%%%%%%%%%%%%%%%%
%%%%%%%%%%%%%%%%%%%%%%%%%%%%%%%%%%%%%%%%%%%%%%%%%%%%%%%%%%%%%%%%%%%%%%%%%%
%% Nota
%%%%%%%%%%%%%%%%%%%%%%%%%%%%%%%%%%%%%%%%%%%%%%%%%%%%%%%%%%%%%%%%%%%%%%%%%%
%% Si ricorda che il numero massimo di facciate e' 30.
%% Nel conteggio delle facciate sono incluse 
%%   indice
%%   sommario
%%   capitoli
%% Dal conteggio delle facciate sono escluse
%%   frontespizio
%%   ringraziamenti
%%   allegati    
%%%%%%%%%%%%%%%%%%%%%%%%%%%%%%%%%%%%%%%%%%%%%%%%%%%%%%%%%%%%%%%%%%%%%%%%%%
%%%%%%%%%%%%%%%%%%%%%%%%%%%%%%%%%%%%%%%%%%%%%%%%%%%%%%%%%%%%%%%%%%%%%%%%%%

    % indice
    \tableofcontents
    \clearpage
    
    
          
    % gruppo per definizone di successione capitoli senza interruzione di pagina
%    \begingroup
      % nessuna interruzione di pagina tra capitoli
      % ridefinizione dei comandi di clear page
%      \renewcommand{\cleardoublepage}{} 
%      \renewcommand{\clearpage}{} 
      % redefinizione del formato del titolo del capitolo
      % da formato
      %   Capitolo X
      %   Titolo capitolo
      % a formato
      %   X   Titolo capitolo
      
%      \titleformat{\chapter}
%        {\normalfont\Huge\bfseries}{\thechapter}{1em}{}
%        
%      \titlespacing*{\chapter}{0pt}{0.59in}{0.02in}
%      \titlespacing*{\section}{0pt}{0.20in}{0.02in}
%      \titlespacing*{\subsection}{0pt}{0.10in}{0.02in}
      
    % sommario
    \chapter*{Abstract} % senza numerazione
\label{ch:abstract}

\addcontentsline{toc}{chapter}{Abstract} % da aggiungere comunque all'indice

%Sommario è un breve riassunto del lavoro svolto dove si descrive l'obiettivo, l'oggetto della tesi, le 
%metodologie e le tecniche usate, i dati elaborati e la spiegazione delle conclusioni alle quali siete arrivati.  
%
%Il sommario dell’elaborato consiste al massimo di 3 pagine e deve contenere le seguenti informazioni:
%\begin{itemize}
%  \item contesto e motivazioni 
%  \item breve riassunto del problema affrontato
%  \item tecniche utilizzate e/o sviluppate
%  \item risultati raggiunti, sottolineando il contributo personale del laureando/a
%\end{itemize}


%This Theses contains the work I've carried out during the internship, part of the Bachelor's Degree in Computer Science curriculum, at SpazioDati Srl, a technology company based in Trento, Italy.
%
%The work pertained the development of a Change Data Capture system.
%
%My work will be based on a need for the company to move some part of their database to a new database, due to load constraints on the former database.
%The transfer of data would be carried out using Kafka streams, linked to changes in the main database.
%
%The new database would serve primarily their sales team, which would use it for usage analytics for existing customers.
%This database will co-exist with the previous one, i.e. it will not replace the parts of it that it's replicating, but it would instead keep being updated with the change events triggering in the main database.


This Theses presents work carried out during my curricular internship, which pertains the design and implementation of a Change Data Capture system.
An interest of the marketing team gave rise to the demand for the presented system, as explained in Chapter \ref{ch:intro}, which in essence solves the problem of capturing data change events from a data source, sometimes aggregating such data, and store them in a separate place for the marketing team to make further use thereof.

An introduction to the scope and rationale for the system is presented in the first chapter.
The several technologies that constitute the subsystems presented are explained in the second chapter, along with their structural composition into the overall system.
Next, several properties of the data at hand are introduced in Chapter \ref{ch:data} and one of them is explored in greater detail in Chapter \ref{ch:timetravel}.
Additionally, some paths for possible future enhancements are outlined in Chapter \ref{ch:future}.
In Appendix \ref{ch:source} the listings of relevant source code are provided.

As the contents of this document pertain the work I've carried out within SpazioDati, all information relevant to their intellectual properties has been anonymized so as to not disclose it.
The methods for such anonymization have been portrayed in \S \ref{sec:non-disclosure}.
Nowhere in this document their customers' data is shown or cited, partly, or in whole.

%%%%%%%%%%%%%%%%%%%%%%%%%%%%%%%%%%%%%%%%%%%%%%%%%%%%%%%%%%%%%%%%%%%%%%%%%%
%%%%%%%%%%%%%%%%%%%%%%%%%%%%%%%%%%%%%%%%%%%%%%%%%%%%%%%%%%%%%%%%%%%%%%%%%%
%% Nota
%%%%%%%%%%%%%%%%%%%%%%%%%%%%%%%%%%%%%%%%%%%%%%%%%%%%%%%%%%%%%%%%%%%%%%%%%%
%% Sommario e' un breve riassunto del lavoro svolto dove si descrive 
%% l’obiettivo, l’oggetto della tesi, le metodologie e 
%% le tecniche usate, i dati elaborati e la spiegazione delle conclusioni 
%% alle quali siete arrivati.
%% Il sommario dell’elaborato consiste al massimo di 3 pagine e deve contenere le seguenti informazioni: 
%%   contesto e motivazioni
%%   breve riassunto del problema affrontato
%%   tecniche utilizzate e/o sviluppate
%%   risultati raggiunti, sottolineando il contributo personale del laureando/a
%%%%%%%%%%%%%%%%%%%%%%%%%%%%%%%%%%%%%%%%%%%%%%%%%%%%%%%%%%%%%%%%%%%%%%%%%%
%%%%%%%%%%%%%%%%%%%%%%%%%%%%%%%%%%%%%%%%%%%%%%%%%%%%%%%%%%%%%%%%%%%%%%%%%%      
      
    %%%%%%%%%%%%%%%%%%%%%%%%%%%%%%%%
    % lista dei capitoli
    %
    % \input oppure \include
    %
    \chapter{Introduction}
\label{ch:intro}


    \chapter{Technologies Employed}
\label{cha:techs}

\section{Apache Kafka}

\subsection{Kafka Connect}

\subsection{Kafka Streams}

\section{PostgreSQL}

\subsection{Write Ahead Log}

    \chapter{Remarks on the Data}
\label{ch:data}

\newcommand{\ind}{\ensuremath{\phantom{= =}}}


We will assume the following, without proof:

\begin{axiom}
	$\source{*}$ is inherently consistent.
\end{axiom}

The consistency of the source database should be assumed because a replication of a database can only be as consistent as the original database, i.e. there is no way to have consistent data, starting from inconsistent data.

Additionally, such is safe to assume, since the database constraints are checked at the source by the PostgreSQL server.

Generally, 
$\dest{\alpha} := \pi_{S}\left(\source{\alpha}\right)$
over some subset $S$ of its columns,
with one exception for table $\gamma$:

$$
\dest{\gamma} := \pi_{\dest{\alpha}.\text{id}, \source{\gamma}.*}
\left(
\source{\gamma} \bowtie_{\alpha.\text{email} = \gamma.\text{email}} \dest{\alpha}
\right)
$$


\section{Aggregations}

\begin{align}
&\dest{\eta} := M\bigl(\sigma^*_{n := \text{count}(*), C := \sum\text{cost}} \bigl( \nonumber \\
&\ind G_{\text{user, time, class}} \bigl(\source{\eta}\bigr) \nonumber \\
&\bigr)\bigr) \nonumber
\end{align}


\section{On Constraints}
\subsection{On Foreign Key Constraints}

    \chapter{Time Travel Functionality}
\label{ch:timetravel}

Every table has one representation in the source database and another representation in the destination database.
To avoid ambiguity, let us introduce the following notation:

\begin{itemize}
\item $\source{\omega}$ and $\dest{\omega}$ refer to the representation of some table $\omega$ at the source and destination databases respectively;
\item $\source{*}$ and $\dest{*}$ refer to the representations of all tables collectively (i.e. the whole databases), at the source and destination databases respectively.
\end{itemize}


We will assume the following, without proof:

\begin{axiom}
$\source{*}$ is inherently consistent.
\end{axiom}

The consistency of the source database should be assumed because a replication of a database can only be as consistent as the original database, i.e. there is no way to have consistent data, starting from inconsistent data.

Additionally, such is safe to assume, since the database constraints are checked at the source by the PostgreSQL server.

Generally, 
$\dest{\alpha} := \pi_{S}\left(\source{\alpha}\right)$
over some subset $S$ of its columns,
with one exception for table $\gamma$:

$$
\dest{\gamma} := \pi_{\dest{\alpha}.\text{id}, \source{\gamma}.*}
    \left(
        \source{\gamma} \bowtie_{\alpha.\text{email} = \gamma.\text{email}} \dest{\alpha}
    \right)
$$

\section{Empty Validity Ranges}

When dealing with ranges, one should account for the fact that $\emptyset$ is indeed a valid range:

$$
\left[t, t\right) = \emptyset, \forall t
$$

This has the inherent meaning, in our domain, that two different records for the same table were issued at the exact same time $t$.

At first, this may seem an issue that could possibly be dismissed as very unlikely.
Nevertheless, the presence of database transactions makes the significance of this scenario become evident.
Within a transaction, different queries are evaluated and later performed as a single atomic\footnote{??? explain atomic} operation, thus resulting in several records being issued at the same instant.


    \chapter{Future Work}
\label{ch:future}


\section{Keeping up to date}
With the source db changes

\section{Metrics}

\section{Batching}

\section{Squashing Transactions Together}

With Kafka Streams

    \chapter{Conclusion}
\label{ch:conclusion}

      
%      \endgroup


    % bibliografia in formato bibtex
    %
    % aggiunta del capitolo nell'indice
%    \addcontentsline{toc}{chapter}{Bibliografia}
    % stile con ordinamento alfabetico in funzione degli autori
    %\bibliographystyle{plain}
    %\bibliography{biblio}\addcontentsline{toc}{chapter}{Bibliography}
%    \nocite{*}
    \printbibliography\addcontentsline{toc}{chapter}{Bibliography}

%%%%%%%%%%%%%%%%%%%%%%%%%%%%%%%%%%%%%%%%%%%%%%%%%%%%%%%%%%%%%%%%%%%%%%%%%%
%%%%%%%%%%%%%%%%%%%%%%%%%%%%%%%%%%%%%%%%%%%%%%%%%%%%%%%%%%%%%%%%%%%%%%%%%%
%% Nota
%%%%%%%%%%%%%%%%%%%%%%%%%%%%%%%%%%%%%%%%%%%%%%%%%%%%%%%%%%%%%%%%%%%%%%%%%%
%% Nella bibliografia devono essere riportati tutte le fonti consultate 
%% per lo svolgimento della tesi. La bibliografia deve essere redatta 
%% in ordine alfabetico sul cognome del primo autore. 
%% 
%% La forma della citazione bibliografica va inserita secondo la fonte utilizzata:
%% 
%% LIBRI
%% Cognome e iniziale del nome autore/autori, la data di edizione, titolo, casa editrice, eventuale numero dell’edizione. 
%% 
%% ARTICOLI DI RIVISTA
%% Cognome e iniziale del nome autore/autori, titolo articolo, titolo rivista, volume, numero, numero di pagine.
%% 
%% ARTICOLI DI CONFERENZA
%% Cognome e iniziale del nome autore/autori (anno), titolo articolo, titolo conferenza, luogo della conferenza (città e paese), date della conferenza, numero di pagine. 
%% 
%% SITOGRAFIA
%% La sitografia contiene un elenco di indirizzi Web consultati e disposti in ordine alfabetico. 
%% E’ necessario:
%%   Copiare la URL (l’indirizzo web) specifica della pagina consultata
%%   Se disponibile, indicare il cognome e nome dell’autore, il titolo ed eventuale sottotitolo del testo
%%   Se disponibile, inserire la data di ultima consultazione della risorsa (gg/mm/aaaa).    
%%%%%%%%%%%%%%%%%%%%%%%%%%%%%%%%%%%%%%%%%%%%%%%%%%%%%%%%%%%%%%%%%%%%%%%%%%
%%%%%%%%%%%%%%%%%%%%%%%%%%%%%%%%%%%%%%%%%%%%%%%%%%%%%%%%%%%%%%%%%%%%%%%%%%
    

%    \titleformat{\chapter}
%        {\normalfont\Huge\bfseries}{Attachment \thechapter}{1em}{}
    % sezione Allegati - opzionale
    \appendix
    \chapter{Listings of Source Files}
\label{ch:source}

A selection of the source code that runs the system presented in this Thesis, is here listed.

In previous chapters, the name of \texttt{class} has been assigned to several relation attributes; note that since \texttt{class} is a reserved keyword in Scala, it will be replaced by \texttt{klass} throughout this Appendix, without implying any change of meaning.


\definecolor{codegreen}{rgb}{0,0.6,0}
\definecolor{codegray}{rgb}{0.5,0.5,0.5}
\definecolor{codepurple}{rgb}{0.58,0,0.82}
\definecolor{backcolour}{rgb}{0.95,0.95,0.92}


\lstdefinestyle{mystyle}{
%    backgroundcolor=\color{backcolour},   
    commentstyle=\color{codegray},
    keywordstyle=\textbf,
    numberstyle=\footnotesize\ttfamily\color{codegray},
%    stringstyle=\color{codepurple},
    basicstyle=\ttfamily,
    breakatwhitespace=false,         
    breaklines=true,
%    captionpos=b,
    keepspaces=true,
    numbers=left,
    numbersep=5pt,
    showspaces=false,
%    showstringspaces=false,
%    showtabs=false,
%    tabsize=2
}

\lstset{style=mystyle}

% File automatically generated by make_source.py
% at 2021-01-25 19:05:46.746438.
% Do not edit.



\lstinputlisting[
    language=Scala,
    caption={source/eu.spaziodati.metrics/streams/StreamsLoader.scala},
    label={src:eu.spaziodati.metrics/streams/StreamsLoader.scala}
]{source/eu.spaziodati.metrics/streams/StreamsLoader.scala}

\lstinputlisting[
    language=Scala,
    caption={source/eu.spaziodati.metrics/streams/etaStream.scala},
    label={src:eu.spaziodati.metrics/streams/etaStream.scala}
]{source/eu.spaziodati.metrics/streams/etaStream.scala}

\lstinputlisting[
    language=Scala,
    caption={source/eu.spaziodati.metrics/streams/kappaStream.scala},
    label={src:eu.spaziodati.metrics/streams/kappaStream.scala}
]{source/eu.spaziodati.metrics/streams/kappaStream.scala}

\lstinputlisting[
    language=Scala,
    caption={source/eu.spaziodati.metrics/streams/lambdaStream.scala},
    label={src:eu.spaziodati.metrics/streams/lambdaStream.scala}
]{source/eu.spaziodati.metrics/streams/lambdaStream.scala}

\lstinputlisting[
    language=Scala,
    caption={source/eu.spaziodati.metrics/streams/thetaStream.scala},
    label={src:eu.spaziodati.metrics/streams/thetaStream.scala}
]{source/eu.spaziodati.metrics/streams/thetaStream.scala}

\lstinputlisting[
    language=Scala,
    caption={source/eu.spaziodati.metrics/connectors/PostgresSinkTask.scala},
    label={src:eu.spaziodati.metrics/connectors/PostgresSinkTask.scala}
]{source/eu.spaziodati.metrics/connectors/PostgresSinkTask.scala}



\end{document}
